\begin{lemma}
\label{thm: cannot exit sliding mode if not all states are better}

Suppose $q$ and $q_+$ satisfy \eqref{eq: definition of q+ greedy section}.
The following statements hold 
% for all action values $q$ and $q_+$ that satisfy \eqref{eq: definition of q+ greedy section},
for every policy $\pol \in \grpol(q)$, $\pol_+ \in \grpol(q_+)$, and state $s \in \sset$ where the update distribution $\update$ satisfies $\update(s, \pol) > \update(s, \pol_+)$:
% is biased towards the action $\pol(s)$:
% For every state $s$ that is updated and where the policy changes,
% i.e. $\update(s, \pol)>0$, $\update_+(s, \pol_+)>0$ and $\pol(s) \neq \pol_+(s)$, 
% i.e. $\update(s, \pol)>0$ and $\pol(s) \neq \pol_+(s)$,
\begin{enumerate}
    \item If $q_{\pol}(s, \pol_+) < q_{\pol}(s, \pol)$, then
    % \item If $\aval_{\pol}(s, \pol_+) < 0$, it follows that
    % it holds that
    % \begin{align}
    % \label{eq: lemma a v a advantage polk}
    %     \aval_{\pol}(s, \pol_+(s)) < 0 ,
    % \end{align}
    \begin{align}
        \nonumber
        &q_{\pol}(s, \pol_+) 
        < q_{\pol}(s, \pol) 
        < q_+(s, \pol) 
        \\
        \label{eq: q approach from above}
        &\hspace{2.5cm} 
        \leq q_+(s, \pol_+)
        \leq q(s, \pol_+) \leq q(s, \pol) .
    \end{align}

    \item If $q_{\pol}(s, \pol_+) < q_{\pol}(s, \pol)$ and $v_{\pol_+}(s) < v_{\pol}(s)$, then 
    % along with \eqref{eq: q approach from above}
    % \item If $\aval_{\pol}(s, \pol_+) < 0$ and $\sval_{\pol}(s) > \sval_{\pol_+}(s)$ it follows that
    % in addition to condition \eqref{eq: lemma a v a advantage polk}, it holds that
    % \begin{align}
    % \label{eq: lemma a v a value}
    %     \sval_{\pol}(s) > \sval_{\pol_+}(s) ,
    % \end{align}
    \begin{align}
    \label{eq: q+ approach from above}
        q_{\pol_+}(s, \pol_+) < q_+(s, \pol) \leq q_+(s, \pol_+) .
    \end{align}

    \item If $q_{\pol}(s, \pol_+) < q_{\pol}(s, \pol)$, $v_{\pol_+}(s) < v_{\pol}(s)$ and $q_{\pol_+}(s, \pol_+) < q_{\pol_+}(s, \pol)$, then
    % along with \eqref{eq: q approach from above} and \eqref{eq: q+ approach from above}
    % \item If $\aval_{\pol}(s, \pol_+) < 0$, $\sval_{\pol}(s) > \sval_{\pol_+}(s)$ and $\aval_{\pol_+}(s, \pol) > 0$, it follows that
    % in addition to conditions \eqref{eq: lemma a v a advantage polk} and \eqref{eq: lemma a v a value}, it holds that
    % \begin{align}
    % \label{eq: lemma a v a advantage polkpo}
    %     \aval_{\pol_+}(s, \pol(s)) > 0 ,
    % \end{align}
    \begin{align}
        \nonumber
        &\update_+(s, \pol) (\qval_{\pol_+}(s, \pol) - \qval_+(s, \pol))
        \\
        & \hspace{2cm} >
        \update_+(s, \pol_+) (\qval_{\pol_+}(s, \pol_+) - \qval_+(s, \pol_+))
    \end{align}
    for every update distribution $\update_+$ where $\update_+(s, \pol_+) > \update_+(s, \pol)$.
\end{enumerate}
\end{lemma}

\begin{proof}
(1) Suppose that in state $s$
\begin{align}
    \label{eq: cond 1 on qpol}
    &q_{\pol}(s, \pol_+) < q_{\pol}(s, \pol) .
\end{align}
Since the policies $\pol$ and $\pol_+$ are greedy with respect to $q$ and $q_+$ we have
\begin{align}
    \label{eq: cond 2 on q}
    &q(s, \pol) \geq q(s, \pol_+) ,
    \\
    \label{eq: cond 3 on q+}
    &q_+(s, \pol) \leq q_+(s, \pol_+) .
\end{align}
By combining \eqref{eq: cond 1 on qpol} with \autoref{thm: all inequalities when switching policy}(2), we obtain that
\begin{align*}
    \update(s, \pol) ( \qval_{\pol}(s, \pol) - \qval(s, \pol_+) )
    &\leq
    \update(s, \pol_+) ( \qval_{\pol}(s, \pol_+) - \qval(s, \pol_+) )
    \\
    &< 
    \update(s, \pol_+) ( \qval_{\pol}(s, \pol) - \qval(s, \pol_+) )
\end{align*}
which simplifies to
\begin{align*}
% \label{eq: first ineq for greedy proof}
    ( \update(s, \pol) - \update(s, \pol_+) )
    ( q_{\pol}(s, \pol) - q(s, \pol_+) ) < 0.
\end{align*}
Therefore, if $\update(s, \pol) > \update(s, \pol_+)$, it follows that
\begin{align*}
% \label{eq: approach from above}
    q_{\pol}(s, \pol)
    <
    q(s, \pol_+) ,
\end{align*}
and together with inequalities \eqref{eq: cond 1 on qpol} and \eqref{eq: cond 2 on q} we see that
\begin{align}
\label{eq: first set of conditions}
    q_{\pol}(s, \pol_+)
    < q_{\pol}(s, \pol)
    < q(s, \pol_+) 
    \leq q(s, \pol) .
\end{align}

Further, the definition of $q_+$ indicates that
\begin{align*}
    q_+(s, \pol)
    &= ( 1 - \lrate \update(s, \pol) ) q(s, \pol) 
    + \lrate \update(s, \pol) q_{\pol}(s, \pol) ,
    \\
    q_+(s, \pol_+)
    &= ( 1 - \lrate \update(s, \pol_+) ) q(s, \pol_+) 
    + \lrate \update(s, \pol_+) q_{\pol}(s, \pol_+) .
\end{align*}
Without loss of generality, we assume that the effective learning rate $\lrate \update$ is positive and smaller than one for all state-action pairs.
The inequalities in \eqref{eq: first set of conditions} then imply that
\begin{align}
\label{eq: condition on s pol}
    &q_{\pol}(s, \pol) < q_+(s, \pol) < q(s, \pol) ,
    \\
\label{eq: condition on s pol+}
    &q_{\pol}(s, \pol_+) < q_+(s, \pol_+) \leq q(s, \pol_+) .
\end{align}
Altogether, expression \eqref{eq: q approach from above} is \eqref{eq: cond 1 on qpol} followed by the first inequality in \eqref{eq: condition on s pol}, then \eqref{eq: cond 3 on q+}, the second inequality in \eqref{eq: condition on s pol+} and finally \eqref{eq: cond 2 on q}.


(2) Suppose that in state $s$
\begin{align*}
    q_{\pol_+}(s, \pol_+) = v_{\pol_+}(s) < v_{\pol}(s) = q_{\pol}(s, \pol) .
\end{align*}
Expression \eqref{eq: q+ approach from above} then follows directly from \eqref{eq: q approach from above}.


(3) Suppose that in state $s$
\begin{align}
    \label{eq: cond on qpol+}
    &q_{\pol_+}(s, \pol_+) < q_{\pol_+}(s, \pol) .
\end{align}
Since policy $\pol_+$ is greedy with respect to $q_+$, it follows that
\begin{align}
    q_{\pol_+}(s, \pol_+) - q_+(s, \pol_+) 
    < q_{\pol_+}(s, \pol) - q_+(s, \pol) .
\end{align}
Expression \eqref{eq: q+ approach from above} indicates that the left-hand side is strictly negative. 
Thus for every update distribution $\update_+$ satisfying $\update_+(s, \pol_+) > \update_+(s, \pol)$, the order of the inequality remains when multiplying both sides as
\begin{multline*}
    \update_+(s, \pol_+) (\qval_{\pol_+}(s, \pol_+) - \qval_{+}(s, \pol_+))
    \\
    < 
    \update_+(s, \pol) ( \qval_{\pol_+}(s, \pol) - \qval_{+}(s, \pol) ) .
\end{multline*}
\end{proof}

% \begin{lemma}
% \label{thm: cannot exit sliding mode if not all states are better}
% Consider the update defined in equation \eqref{eq: definition of q+ greedy section}, and suppose the update map $\updateset$ satisfies \autoref{asp: on-policy bias}.
% For any state $s$ that is updated and where the policy changes,
% i.e. $\update(s, \pol)>0$ and $\pol(s) \neq \pol_+(s)$, 
% if
% \begin{align}
%     \label{eq: lemma a v a advantage polk}
%     &\aval_{\pol}(s, \pol_+(s)) < 0 ,
%     \\
%     \label{eq: lemma a v a value}
%     &\sval_{\pol}(s) > \sval_{\pol_+}(s) ,
%     \\
%     \label{eq: lemma a v a advantage polkpo}
%     &\aval_{\pol_+}(s, \pol(s)) > 0 ,
% \end{align}
% then
% \begin{multline}
% \label{eq: switch inequality on side pol_new}
%     \update_+(s, \pol) (\qval_{\pol}(s, \pol) - \qval_+(s, \pol))
%     \\
%     >
%     \update_+(s, \pol_+) (\qval_{\pol_+}(s, \pol_+) - \qval_+(s, \pol_+)) .
% \end{multline}
% % Moreover, if policy $\pol_\kpo$ is strictly greedy on $q_\kpo$, i.e. $\qval_{\kpo}(s, \pol_\kpo) > \qval_{\kpo}(s, \pol_\kpz)$, then inequality \eqref{eq: switch inequality on side pol_new} is strict.
% \end{lemma}

% \begin{proof}
% % For any state $s$ where $\pol_k(s) \neq \pol_\kpo(s)$, 
% Conditions \eqref{eq: lemma a v a advantage polk} and \eqref{eq: lemma a v a value} imply, by \autoref{thm: approach from above}, that
% \begin{align*}
% % \label{eq: approach from above kpo}
%     \qval_{\kpo}(s, \pol_\kpz) > \qval_{\pol_\kpo}(s, \pol_\kpo) .
%     % \leq \qval_{\kpo}(s, \pol_\kpo) .
%     % 0 >
%     % \qval_{\pol_\kpo}(s, \pol_\kpo) - \qval_{\kpo}(s, \pol_\kpz)
%     % \geq \qval_{\pol_\kpo}(s, \pol_\kpo) - \qval_{\kpo}(s, \pol_\kpo) .
% \end{align*}
% Thus, since the policy $\pol_\kpo$ is greedy on $q_\kpo$, i.e. $\qval_{\kpo}(s, \pol_\kpz) \leq \qval_{\kpo}(s, \pol_\kpo)$, it follows that
% \begin{align*}
% % \label{eq: greedy policy}
%     \qval_{\pol_\kpo}(s, \pol_\kpo) - \qval_{\kpo}(s, \pol_\kpo) 
%     \leq \qval_{\pol_\kpo}(s, \pol_\kpo) - \qval_{\kpo}(s, \pol_\kpz)
%     < 0 .
% \end{align*}
% As a result, given that the difference on the left-hand side is strictly negative and that, by \autoref{asp: on-policy bias}, $\update_\kpo(s, \pol_\kpz(s)) \leq \update_\kpo(s, \pol_\kpo(s))$, we observe that
% \begin{align*}
%     \update_\kpo(s, \pol_\kpo) (\qval_{\pol_\kpo}(s, \pol_\kpo) - \qval_{\kpo}(s, \pol_\kpo))
%     &\leq 
%     \update_\kpo(s, \pol_\kpz) ( \qval_{\pol_\kpo}(s, \pol_\kpo) - \qval_{\kpo}(s, \pol_\kpo) )
%     \\
%     &\leq 
%     \update_\kpo(s, \pol_\kpz) ( \qval_{\pol_\kpo}(s, \pol_\kpo) - \qval_{\kpo}(s, \pol_\kpz) )
%     \\
%     &< 
%     \update_\kpo(s, \pol_\kpz) ( \qval_{\pol_\kpo}(s, \pol_\kpz) - \qval_{\kpo}(s, \pol_\kpz) ) ,
% \end{align*}
% where we used condition \eqref{eq: lemma a v a advantage polkpo} in the last step.
% % The last inequality rewrites to
% % \begin{multline*}
% %     \update_\kpo(s, \pol_\kpz) (\qval_{\pol_\kpo}(s, \pol_\kpz) - \qval_\kpo(s, \pol_\kpz))
% %     \\
% %     >
% %     \update_\kpo(s, \pol_\kpo) (\qval_{\pol_\kpo}(s, \pol_\kpo) - \qval_\kpo(s, \pol_\kpo)) .
% % \end{multline*}
% \end{proof}